% ============================================================================
%  back/tests.tex — بانکِ تست‌هایِ اضافی (طبقه‌بندی‌شده بر اساسِ درس و مبحث)
% ============================================================================
\Jpart{بانکِ تست‌هایِ اضافی}
  {تست‌هایِ طبقه‌بندی‌شده بر اساسِ درس و مبحث}
  {تست‌های چهارگزینه‌ایِ فراتر از تمرین‌های کتاب؛ پاسخنامهٔ کاملاً تشریحی در
   انتهای همین بخش}

% ============================================================================
%  بخشِ اوّل: درسِ اوّل
% ============================================================================
\Jsec{بخشِ ١ \,\text| تست‌هایِ درسِ اوّل}

\Jsub{مبحث: حروفِ مشبههٔ فعل}

\JQ{در کدام گزینه حرفِ مشبههٔ فعل، «آرزو» را می‌رساند؟}
{لَعَلَّكُمْ تَعْقِلُونَ}{وَلٰكِنَّ أَكْثَرَ النَّاسِ لا يَشْكُرُونَ}
{يَا لَيْتَنِي كُنْتُ تُرَابًا}{إِنَّ اللّٰهَ غَفُورٌ رَحِيمٌ}

\JQ{اعرابِ کدام واژه در جملهٔ «إِنَّ الْمُؤْمِنِينَ إِخْوَةٌ» درست است؟}
{الْمُؤْمِنِينَ: خبرِ إِنَّ}{إِخْوَةٌ: خبرِ إِنَّ، مرفوع}
{إِخْوَةٌ: اسمِ إِنَّ، منصوب}{الْمُؤْمِنِينَ: مبتدا، مرفوع}

\JQ{«پدرم گفت که علم کلیدِ زندگی است» — کدام گزینه ترجمهٔ درستِ بخشِ «که» است؟}
{إِنَّ}{أَنَّ}{لٰكِنَّ}{لَعَلَّ}

\JQ{در کدام گزینه، اسم و خبرِ حرفِ مشبهه هر دو «مرفوع»اند؟}
{إِنَّ الْعِلْمَ نَافِعٌ}{كَأَنَّ الطَّالِبَةَ مُجِدَّةٌ}
{لَعَلَّ الْجَوَّ صَافٍ}{أَنَّ الرَّجُلَ صَادِقٌ... هر دو مرفوع‌اند}

\JQ{«لِمَاذَا نَجَحْتَ؟» — پاسخِ مناسب با ساختارِ درست:}
{لِأَنَّ أَنَا مُجْتَهِدٌ}{لِأَنِّي مُجْتَهِدٌ}
{إِنِّي مُجْتَهِدٌ}{كَأَنِّي مُجْتَهِدٌ}

\Jsub{مبحث: لا نافیهٔ جنس}

\JQ{در جملهٔ «لا حَوْلَ وَ لا قُوَّةَ إِلَّا بِاللّٰهِ»، «حَوْلَ» چه نقشی دارد؟}
{مفعول}{اسمِ لا نافیهٔ جنس}{مضافٌ‌الیه}{حال}

\JQ{کدام گزینه «لا نافیهٔ جنس» دارد؟}
{لا تَذْهَبْ إِلَى السُّوقِ}{لا، أَنَا مِنْ طَبْرِيزَ}
{لا يَعْلَمُ الْغَيْبَ إِلَّا اللّٰهُ}{لا أَحَدَ فِي الْفَصْلِ}

\JQ{«لا شَكَّ فِيهِ» — دربارهٔ ساختارِ این عبارت کدام گزینه درست است؟}
{اسمِ لا: شَكَّ ؛ خبر: فِيهِ}{اسمِ لا: شَكَّ ؛ خبر: حذف شده}
{شَكَّ: خبرِ لا ؛ خبرِ واقعی: فِيهِ}{شَكَّ: مفعولِ حذف‌شده}

\JQ{در «لا نافیهٔ جنس»، اسمِ لا معمولاً:}
{با «ال» و مفتوح است}{با «ال» و مرفوع است}{بدونِ «ال» و مفتوح است}{بدونِ «ال» و مجرور است}

\Jsub{مبحث: ترجمه و واژگانِ درسِ اوّل}

\JQ{معنای واژهٔ «اَلدَّاء» در ابیاتِ درس کدام است؟}
{دعا}{درمان، دارو}{بیماری}{طبیب}

\JQ{«فَفُزْ بِعِلْمٍ وَ لا تَطْلُبْ بِهِ بَدَلاً» — معنای «بَدَل» در این بیت:}
{جانشین، بهرهٔ مادی}{قهرمان}{پیکر}{گِل}

\JQ{مفردِ واژهٔ «عِظَام» کدام است؟}
{عَظْم}{عَظِيم}{عَظَمَة}{أَعَاظِم}

% ============================================================================
%  بخشِ دوم: درسِ دوم
% ============================================================================
\Jsec{بخشِ ٢ \,\text| تست‌هایِ درسِ دوم}

\Jsub{مبحث: اَلْحَال (قیدِ حالت)}

\JQ{«حال» در جملهٔ «وَصَلَ الْمُسَافِرُونَ مُتَأَخِّرِينَ» کدام است و چه
حرکتی دارد؟}
{الْمُسَافِرُونَ — مرفوع}{مُتَأَخِّرِينَ — منصوب}
{مُتَأَخِّرِينَ — مرفوع}{الْمُطَار — مجرور}

\JQ{در کدام گزینه، واژهٔ زیرخط‌دار «صفت» است نه «حال»؟}
{رَأَيْتُ رَجُلًا مَسْرُورًا}{رَأَيْتُ الْرَّجُلَ الْمَسْرُورَ}
{خَرَجَ الطَّالِبُ مُبْتَسِمًا}{رَجَعَ الْجُنُودُ فَرِحِينَ}

\JQ{ترجمهٔ درستِ «نَظَرْتُ إِلَيْهِ وَ هُوَ يَكْتُبُ» کدام است؟}
{به او نگاه کردم و نوشت}{به او نگاه کردم، در حالی‌که می‌نوشت}
{به او نگاه کردم در حالی‌که نوشت}{نگاهِ او به نوشتنِ من بود}

\JQ{واو در جملهٔ «دَخَلَ الْمُدِيرُ وَ هُوَ مُبْتَسِمٌ» چه نوعی است؟}
{واوِ عطف}{واوِ حالیّه}{واوِ معیه}{واوِ قسم}

\Jsub{مبحث: تحلیلِ صرفی، اعراب و واژگان}

\JQ{«اَلْعُمَّالُ» در جملهٔ «اَلْعُمَّالُ يَشْتَغِلُونَ» چه نوعی است؟}
{اسمِ مکان}{جمعِ تکسیرِ اسمِ فاعل}{اسمِ مبالغه}{اسمِ مفعول}

\JQ{مفردِ «اَلْمَنَاجِم» کدام است؟}
{مَنْجَم}{أَنْجُم}{نُجُوم}{مَنَاجِمَة}

\JQ{کدام گزینه «جمعِ مکسّر» دارد؟}
{اَلْمَجَالَات}{اَلصِّبْيَان}{اَلْأَنْفَاق}{اَلْأَنْحَاء}

\JQ{واژهٔ «اَلدَّؤُوب» به معنای ... است.}
{سریع‌الانفجار}{با پشتکار}{ویران‌شده}{ثروتمند}

\JQ{«وَ إِنْ كَانَ غَرَضُهُ مِنِ اخْتِرَاعِهِ مُسَاعَدَةَ الْإِنْسَانِ» — «غَرَضُهُ»
چه نقشی دارد؟}
{فاعلِ «كَانَ»}{اسمِ «كَانَ» (مبتدا)}{خبرِ «كَانَ»}{مفعول}

\JQ{ترجمهٔ درستِ «هَلْ تُعْطَى الْجَوَائِزُ الْيَوْمَ لِمَنْ هُوَ أَهْلٌ
لِذٰلِكَ؟» کدام است؟}
{آیا جوایز را امروز به اهلش می‌دهند؟}{آیا امروز جوایز به شایستگان اهدا می‌شود؟}
{آیا جوایز امروز داده شد؟}{آیا اهلِ جایزه امروز می‌آیند؟}

% ============================================================================
%  پاسخنامهٔ تشریحی
% ============================================================================
\Jsec{پاسخ‌نامهٔ تشریحیِ بانکِ تست‌ها}

\begin{JKey}[پاسخ‌نامهٔ بخشِ ١ (درسِ اوّل)]
\JAnswer{تستِ ١}{گزینهٔ ۳ — «لَيْتَ» حرفِ آرزوست («یاکاش»...)؛ «لَعَلَّ» امید،
«لٰكِنَّ» استثنا/تقابل و «إِنَّ» تأکید می‌رساند.}
\JAnswer{تستِ ٢}{گزینهٔ ۲ — «الْمُؤْمِنِينَ» اسمِ إِنَّ و منصوب به یایِ جمع
(اِسْمِه) است و «إِخْوَةٌ» خبرِ مرفوع به ضمّه. گزینه‌های دیگر اعراب را جابه‌جا
گفته‌اند.}
\JAnswer{تستِ ٣}{گزینهٔ ۲ — «أَنَّ» معنای «که» دارد و نقل‌قولِ غیرمستقیم
(گفتِ پدر) را به جملهٔ بعد پیوند می‌دهد؛ «إِنَّ» برای تأکیدِ جملهٔ مستقل است.}
\JAnswer{تستِ ٤}{گزینهٔ ۲ — «كَأَنَّ الطَّالِبَةَ مُجِدَّةٌ»: اسمِ كَأَنَّ
(الطَّالِبَةَ) منصوب و خبرِ آن (مُجِدَّةٌ) مرفوع است؛ گزینهٔ ۱ و ۳ خبر را
منصوب بسته‌اند که غلط است.}
\JAnswer{تستِ ٥}{گزینهٔ ۲ — «لِأَنِّي» ترکیبِ لِأَنَّ + ضمیرِ متکلمِ منفصلِ
«إِيَّا» به‌صورتِ «لِأَنِّي» است (چسبیده). پاسخِ کاملِ «چرا موفق شدی؟»:
«زیرا من کوشا هستم».}
\JAnswer{تستِ ٦}{گزینهٔ ۲ — «حَوْلَ» اسمِ لا نافیهٔ جنس است (بدونِ «ال» و
مفتوح)؛ معنا: «هیچ جنبش‌وقدرتی نیست مگر به خدا».}
\JAnswer{تستِ ٧}{گزینهٔ ۴ — «لا أَحَدَ فِي الْفَصْلِ»: هیچ‌کس در کلاس نیست
(اسمِ لا: أَحَدَ). گزینهٔ ۱ لای نهی، ۲ لا در پاسخِ پرسش و ۳ لا نفیِ مضارع است.}
\JAnswer{تستِ ٨}{گزینهٔ ۲ — در «لا شَكَّ فِيهِ» خبرِ لا حذف شده (لا شَكَّ
[مَوْجُودٌ] فِيهِ)؛ «فِيهِ» جارّ و مجرور است، نه خودِ خبرِ ظاهر.}
\JAnswer{تستِ ٩}{گزینهٔ ۳ — اسمِ لا نافیهٔ جنس معمولاً بدونِ «ال» و دارایِ
فتحه است؛ مجرور بودن از نشانه‌های اسمِ لا نیست.}
\JAnswer{تستِ ١٠}{گزینهٔ ۳ — «دَاء» = بیماری (جمعِ باب‌های اشتهایی: أَدْوَاء)؛
دامِ آزمون، معنای «دارو» است که معنای واژهٔ «دَاوَاة/دُوَاء» است نه «دَاء».
در بیتِ درس: «دَوَاؤُكَ فِيكَ» یعنی بیماری‌ات در توست.}
\JAnswer{تستِ ١١}{گزینهٔ ۱ — «بَدَل» یعنی جانشین/بهرهٔ مادی؛ مفهومِ بیت: علم
را فدای مادیات مکن.}
\JAnswer{تستِ ١٢}{گزینهٔ ۱ — عِظَام جمعِ مکسّرِ «عَظْم» (استخوان) است؛
«أَعَاظِم» جمعِ «عَظِيم/عَظْمَة» است.}
\end{JKey}

\begin{JKey}[پاسخ‌نامهٔ بخشِ ٢ (درسِ دوم)]
\JAnswer{تستِ ١}{گزینهٔ ۲ — حال همیشه منصوب است؛ «مُتَأَخِّرِينَ» جمعِ مذکرِ
سالمِ منصوب (بیه‌بعد از یایِ جمع حذف می‌شود) و حالِ «الْمُسَافِرُونَ» است.}
\JAnswer{تستِ ٢}{گزینهٔ ۲ — «الْمَسْرُورَ» با «ال» و مطابقِ موصوف، صفتِ
«الرَّجُلَ» است؛ بقیه نکره‌اند و حال‌اند.}
\JAnswer{تستِ ٣}{گزینهٔ ۲ — جملهٔ حالیّه با واوِ حالیّه + ضمیر، هنگامِ فعلِ
مضارع در جملهٔ اصلی، به ماضیِ استمراری ترجمه می‌شود: «در حالی‌که می‌نوشت».}
\JAnswer{تستِ ٤}{گزینهٔ ۲ — واوِ حالیّه است؛ جملهٔ «وَ هُوَ مُبْتَسِمٌ»
جملهٔ حالیّه برای «الْمُدِيرُ» است، نه عطف.}
\JAnswer{تستِ ٥}{گزینهٔ ۲ — «عُمَّال» جمعِ تکسیرِ «عَامِل» (اسمِ فاعلِ «عَمِلَ»)
است؛ اسمِ مکانِ «عَمِلَ» یعنی «مَعْمَل» (کارخانه).}
\JAnswer{تستِ ٦}{گزینهٔ ۱ — مفردِ الْمَنَاجِم یعنی الْمَنْجَم؛ «أَنْجُم» و
«نُجُوم» مربوط به «نَجْم» (ستاره) است.}
\JAnswer{تستِ ٧}{گزینهٔ ۴ — «أَنْحَاء» جمعِ مکسّرِ «نَحْو» است (وزنِ
أَفْعَال)؛ مَجَالَات و أَنْفَاق جمعِ سالم و صِبْيَان جمعِ مؤنثِ سالم‌اند.}
\JAnswer{تستِ ٨}{گزینهٔ ۲ — «الدَّؤُوب» صفتِ مشبهه به معنای «با پشتکار» است.}
\JAnswer{تستِ ٩}{گزینهٔ ۲ — «كَانَ» فعلِ ناقصه است؛ اسمِ آن «غَرَضُهُ»
(مبتدا) و خبرش «مُسَاعَدَةَ الْإِنْسَانِ» (منصوب) است؛ «مِنِ اخْتِرَاعِهِ»
جارّ و مجرور.}
\JAnswer{تستِ ١٠}{گزینهٔ ۲ — فعلِ «تُعْطَى» مجهول است؛ ترجمهٔ دقیق: «آیا
امروز جوایز به کسی که شایستگیِ آن را دارد، داده می‌شود؟».}
\end{JKey}
